\chapter{Related work}\label{chapter:related_work}

  \section{Music Genre classification with audio features}

  \section{EEG datasets and studies with music stimulus}
    As music is a popular stimulus used in \ac{EEG} studies to collect the signal data, many research projects have been conducted in this area. This section intends to serve as an overview of work that has been done in the context of constructing an EEG dataset where music served as a stimulus. For each dataset, the experiment setup and acquisition methods are described, and certain preprocessing and denoising steps are emphasised.
    \subsection{DEAP}
      The \ac{DEAP} is an \ac{EEG} dataset made public in 2012 by Koelstra et al. It is a multimodal dataset providing data for analysing human affective states. For 32 participants, \ac{EEG} and peripheral physiological signals were recorded while watching 40 one-minute sections of music videos. The video was also shown alongside the music, serving as an additional stimulus. After each snippet, test subjects were asked to rate each music video in terms of valence (positively or negatively connoted), arousal, liking, dominance and familiarity. \cite{koelstra2012deap}\par
      The music comes from a website called Last.fm and was selected using a sophisticated procedure. 120 music videos were initially selected and labelled with keywords out of a list of emotional keywords provided by the Last.fm website. Based on each emotional keyword, ten songs labelled most with one of them have been selected in the next step of the procedure. The chosen songs were further filtered according to valence and arousal attributes. After this filtering step, a one-minute highlight section was selected with an effective highlighting algorithm from each music video. Last, a web-based assessment experiment was conducted to decide the 40 final music video clips for the \ac{DEAP} study. \cite{koelstra2012deap} \par
      The experiment was conducted in a laboratory environment with controlled illumination. The EEG signals have been recorded with a sampling rate of 512Hz by 32 active AgCl electrodes placed in the international 10-20 system. Apart from that, not only have physiological signals been measured, but frontal face videos for 22 of the 32 participants have also been recorded. The basic procedure of each session includes a five-second baseline recording before displaying the one-minute music video. When the video ended, the participant assessed the clip based on arousal, valence, liking, familiarity and dominance. \cite{koelstra2012deap}\par
      For further investigation of the recorded signal data, the researchers provide a preprocessed variant of the \ac{EEG} data. First and foremost, the data were commonly referenced and then downsampled to 256Hz. The downsampled signal went through a high-pass filter with a 2 Hz cutoff frequency employing the EEGlab toolbox. Eye artefacts were removed by using a \ac{BSS} technique. For further analysis, the baseline recording of five seconds was removed. \cite{koelstra2012deap}\par

    \subsection{Music BCI}
      Behavioural responses like in the \ac{DEAP} dataset can be researched, as well as the perception of musical patterns. Therefore, researchers from Berlin and Copenhagen have chosen polyphonic music as a stimulus for an \ac{EEG} dataset. Polyphonic music describes music with two or more parallel melodic lines, including different instruments. In their work, the main goal is to explore the suitability of polyphonic music in \ac{BCI}. \cite{treder2014decoding}\par
      For this purpose, eleven participants aged 21-50 years were chosen, each not educated in the field of \ac{BCI} research. For acquisition, a 1000 Hz sampling rate was determined with 64 electrodes placed according to the international 10-10 system. The stimuli were 40-second music clips, each of which contained six tracks. Two of these tracks are meant for stereo audio playback, the other four are a mono mix of the stereo clip recorded as additional \ac{EEG} channels. Each clip contains three instruments playing musical patterns, where one instrument plays a deviant version while the other two stay the same. Per instrument, there are 3-7 deviations per clip. Synth-pop, which has around 130 bpm, and jazz, which has around 120 bpm, were chosen as genres. While the Synth-Pop clips consist of drums, a bass, and a keyboard, the clips coming from Jazz comprise flute, piano, and a double bass. Ten different clips were made for both music types with those kinds of deviant patterns. \cite{treder2014decoding}\par
      During the experiment, the participant was seated comfortably in a chair and instructed to stay still and relaxed, especially to minimise eye movements during the trial. The main experiment is divided into ten blocks with 21 music clips played randomly. The listener was told to pay attention to each of the three instruments in seven different music clips. To further ensure attention, the participant has the task of counting the number of deviants in the cued instrument. Before the following clip, the participant must enter the recognised number of deviations. \cite{treder2014decoding} \par
      The signal data were preprocessed to further analyse the recorded \ac{EEG} data. After downsampling them to 250 Hz, they were low-pass filtered using a Chebyshev filter accompanied by a passband of 42 Hz and a stopband of 49 Hz. The data is sectioned into epochs from -200 ms to 1200 ms for additional denoising. While the prestimulus interval was used for baseline correction, a min-max criterion was employed to discard physiological artefacts. \cite{treder2014decoding} \par
        
    \subsection{MUSIN-G}
      In 2022, researchers published an open-source \ac{EEG} dataset examining the correlation between the brain and music genres. In the experiment, 20 participants listened to naturalistic songs (meaning mostly without lyrics) across 12 genres, ranging from Indian folk music to Goth Rock or Western electronic music. Apart from the genre, subjective characteristics like familiarity and liking have been considered \cite{miyapuram2022Musing}\par
      It is essential to mention that all participants have little musical training and are therefore considered non-musicians. The participant's \ac{EEG} data were recorded by a 128-channel high-density EEG cap while sitting comfortably in a dark environment. The main procedure can be divided into two parts. At the start of the acquisition, a baseline recording was conducted during absolute silence lasting for two minutes. During this process, the subject's eyes were closed. The second part then consists of listening to the songs, where the order is randomised depending on the person hearing the audio. A beep sound instructed the participants to close their eyes when the song started to play. After each song, a ten-second gap was chosen, followed by another beep sound indicating to open their eyes for answering two questions. The participant has to rate the song on a scale of one to five regarding familiarity and liking. After answering the questions, the next song was played with a trailing ten-second gap in between. \cite{miyapuram2022Musing}\par
      Since the data is provided as raw and segmented \ac{EEG} data, the signals can be preprocessed for further analysis. The authors suggest a workflow that includes the EEGLAB toolbox in MATLAB. First, the signals are high-pass linear FIR filtered at 0.2 Hz. This is followed by the CleanLine method, which removes 50Hz of line noise (comparable to a notch filter). Before segmenting the data into each song, the filtered data were down-sampled to 250Hz. By employing a spectrum criterion, bad channels were removed to prepare them for the artefact removal. The \ac{MARA} was used to remove the artefacts, an \ac{ML}-based method for artefact detection and removal. Lastly, removed channels are interpolated with the spherical function of EEGLAB before the average reference is applied to the whole dataset. \cite{miyapuram2022Musing} 
  \section{Emotion classification with EEG signals}

  \section{Music Genre classification with EEG signals}
  